% Source r6 block, retained for companion integration. Not a standalone document.
\section{The centred integer state}\label{sec:state}

The three maps defined next carry out the two steps announced at the end of
Section~\ref{sec:problem}.  The first two clear denominators along the
sequence, so that a rational reciprocal sum leaves a sequence of integers; the
third subtracts from that integer the value it takes on Sylvester's sequence,
leaving one number whose vanishing is the whole question.  All three are maps
between integer tuples and carry no hypothesis of any kind.

For $a,D,C\in\Z$ define
\[
 \den(a,D)=aD,
 \quad
 \tail(a,D,C)=aC-D,
 \quad
 \ctr(a,D,C)=D-(a-1)C ,
\]
the
\lword{Erdos243/ReciprocalTailRigidity.lean}{41}{nextDenState}{denominator
update}, the
\lword{Erdos243/ReciprocalTailRigidity.lean}{45}{nextTailState}{tail update},
and the
\lword{Erdos243/ReciprocalTailRigidity.lean}{49}{centeredState}{centred state}.
Given a sequence $(a_n)$ of integers, whose terms we call the
\emph{multipliers}, we write $D_{n+1}=\den(a_n,D_n)$ for the \emph{denominator
state}, $C_{n+1}=\tail(a_n,D_n,C_n)$ for the \emph{tail state}, and
$E_n=\ctr(a_n,D_n,C_n)$ for the \emph{centred state}, also called the
\emph{error}.  The word \emph{centred} records what the third map does: it
measures $D_n$ against $(a_n-1)C_n$, and $D_n=(a_n-1)C_n$ holds at every index
on Sylvester's sequence (Example~\ref{ex:sylvester}), so $E$ is a deviation
from that sequence rather than a size.  At an index where $E_n<0$ we write
$e_n=-E_n>0$ and call it the \emph{magnitude} of the error there.

\begin{proposition}[update law]\label{res:update}
$\tail(a,D,C)=C-\ctr(a,D,C)$; that is, $C_{n+1}=C_n-E_n$.
\end{proposition}

\begin{proof}
$C-\bigl(D-(a-1)C\bigr)=aC-D$.
\end{proof}

\noindent Formalised as the
\lword{Erdos243/ReciprocalTailRigidity.lean}{57}{nextTailState_eq_sub_centered}{update
law}.  The state therefore moves by exactly the error at each step, which is
what makes the sign of $E$ decisive: where $E$ is positive the state falls,
where $E$ is negative it rises, and where $E$ vanishes it is stationary.

\begin{example}[the Sylvester orbit]\label{ex:sylvester}
Take $a_n=2,3,7,43,1807,\ldots$ with $a_{n+1}=a_n^{2}-a_n+1$, and start the
state at $D_0=C_0=1$.  The two updates give

\begin{center}
\begin{tabular}{@{}rrrrr@{}}
\toprule
$n$ & $a_n$ & $D_n$ & $C_n$ & $E_n$\\
\midrule
$0$ & $2$ & $1$ & $1$ & $0$\\
$1$ & $3$ & $2$ & $1$ & $0$\\
$2$ & $7$ & $6$ & $1$ & $0$\\
$3$ & $43$ & $42$ & $1$ & $0$\\
$4$ & $1807$ & $1806$ & $1$ & $0$\\
\bottomrule
\end{tabular}
\end{center}

\noindent and $D_n=a_n-1$ with $C_n=1$ at every index: if $D_n=a_n-1$ and
$C_n=1$ then $C_{n+1}=a_n-(a_n-1)=1$ and $D_{n+1}=a_n(a_n-1)=a_{n+1}-1$.
Hence $E_n=D_n-(a_n-1)C_n=0$ throughout and the tail state never moves, which
is Proposition~\ref{res:update} in the stationary case.
\end{example}

\paragraph{The intended reading, which is not formalised.}
Let $T_n=\sum_{k\ge n}1/a_k$ and suppose $T_0\in\Q$.  Put $D_0$ equal to a
common denominator and $D_n=D_0a_0\cdots a_{n-1}$, and set $C_n=D_nT_n$.  Then
$C_n\in\Z$ for every $n$, and from $T_n=1/a_n+T_{n+1}$ one gets
$C_{n+1}=a_nC_n-D_n$, which is the tail update.  On Sylvester's sequence
$T_n=1/(a_n-1)$ exactly, so $D_n=(a_n-1)C_n$ and $E_n=0$: the centred state
measures deviation from the Sylvester tail identity, and it vanishes
identically on the Sylvester orbit.

None of this paragraph is formalised.  The Lean module contains no series, no
rationality hypothesis, and no growth hypothesis; it proves identities and
implications about the integer system above, and about its natural-number
realisation, for which the bridging identities are the
\lword{Erdos243/ReciprocalTailRigidity.lean}{1863}{natTail_eq_nextTailState}{tail
realisation}, the
\lword{Erdos243/ReciprocalTailRigidity.lean}{1875}{natDen_eq_nextDenState}{denominator
realisation}, and the
\lword{Erdos243/ReciprocalTailRigidity.lean}{1972}{natTail_eq_sub_centeredState}{signed
update law}.  The identification with reciprocal tails motivates the
definitions.  The formal development does not check that identification.

\begin{proposition}[scale equivariance]\label{res:scale}
For every $s,a,D,C\in\Z$,
\[
 \den(a,sD)=s\,\den(a,D),\qquad
 \tail(a,sD,sC)=s\,\tail(a,D,C),
\]
\[
 \ctr(a,sD,sC)=s\,\ctr(a,D,C).
\]
\end{proposition}

\noindent Formalised as the
\lword{Erdos243/ReciprocalTailRigidity.lean}{65}{nextDenState_scale}{denominator
equivariance}, the
\lword{Erdos243/ReciprocalTailRigidity.lean}{70}{nextTailState_scale}{tail
equivariance}, and the
\lword{Erdos243/ReciprocalTailRigidity.lean}{75}{centeredState_scale}{centred
equivariance}.  The system depends only on the ray through $(D,C)$.  This is
used twice below: once to normalise the finite search of
Appendix~\ref{app:residue}, and once, in Theorem~\ref{res:periodic}, as the
induction step that removes common scale.

\section{The defect identity and its two consequences}\label{sec:defect}

Section~\ref{sec:state} connects the Sylvester recurrence to the error in one
direction only: on Sylvester's sequence the error vanishes.  The converse is
what the problem needs, and it comes from a single identity, which expresses
the deviation of $a_{n+1}$
from the Sylvester successor, multiplied by the tail state $C_{n+1}$, as a
combination of the errors at $n$ and at $n+1$, the first weighted by
$a_n^{2}$.  Its two consequences are
Theorem~\ref{res:eventual}, in which an eventually vanishing error forces the
recurrence provided the tail state is eventually nonzero, and
Theorem~\ref{res:absorb}, in which, under strict centring (the condition
$|E_n|<C_n$ at every index), one vanishing error propagates to every later
index.

Write $\defect(a,a')=a'-\sylnext(a)$ for the deviation of the next term from
the Sylvester successor, the
\lword{Erdos243/ReciprocalTailRigidity.lean}{53}{sylvesterDefect}{Sylvester
defect}.

\begin{theorem}[defect identity]\label{res:defect}
For all $a,a',D,C\in\Z$,
\[
 \defect(a,a')\cdot\tail(a,D,C)
 =a^{2}\,\ctr(a,D,C)-\ctr\bigl(a',\den(a,D),\tail(a,D,C)\bigr),
\]
that is, $\Delta_n\,C_{n+1}=a_n^{2}E_n-E_{n+1}$.
\end{theorem}

\begin{proof}
A direct expansion: both sides equal
$a'aC-a'D-a^{3}C+a^{2}D+a^{2}C-aD-aC+D$.
\end{proof}

\noindent Formalised as the
\lword{Erdos243/ReciprocalTailRigidity.lean}{1769}{sylvesterDefect_mul_nextTailState}{defect
identity}.  It is a polynomial identity in four indeterminates, with no
hypothesis of any kind, and it carries the whole argument twice over: once for
rigidity, and once for absorption.

\begin{example}[one defect and the error it creates]\label{ex:defect}
Continue Example~\ref{ex:sylvester} but replace $a_3=43$ by $a_3=44$.  The
states $D_3=42$ and $C_3=1$ are unchanged, since they depend only on
$a_0,a_1,a_2$, and $E_2=0$ still.  The defect is
$\Delta_2=a_3-\sylnext(a_2)=44-43=1$, and
$E_3=D_3-(a_3-1)C_3=42-43=-1$, so the identity reads
\[
 \Delta_2C_3=1\cdot1=1=49\cdot0-(-1)=a_2^{2}E_2-E_3 .
\]
By Proposition~\ref{res:update} the tail state then rises, $C_4=C_3-E_3=2$.  A
single unit of defect at one index has produced a negative error of magnitude
$1$ at the next, and the state has started to climb.
\end{example}

\begin{theorem}[two vanishing errors force the step]\label{res:step}
Let $a,a',D,C\in\Z$ with $\tail(a,D,C)\ne0$.  If
$\ctr(a,D,C)=0$ and $\ctr\bigl(a',\den(a,D),\tail(a,D,C)\bigr)=0$, then
$a'=\sylnext(a)$.
\end{theorem}

\begin{proof}
By Theorem~\ref{res:defect} the product $\defect(a,a')\cdot\tail(a,D,C)$ is
zero, and the second factor is nonzero by hypothesis.
\end{proof}

\noindent Formalised as the
\lword{Erdos243/ReciprocalTailRigidity.lean}{1781}{sylvesterNext_eq_of_centered_zero}{local
rigidity step}.  The hypothesis $C_{n+1}\ne0$ is not removable: at $C_{n+1}=0$
the identity gives no information about $a'$.

\begin{theorem}[sequence form]\label{res:eventual}
Let $a,D,C:\N\to\Z$ satisfy $D_{n+1}=\den(a_n,D_n)$ and
$C_{n+1}=\tail(a_n,D_n,C_n)$.  If $E_n=0$ for all sufficiently large $n$ and
$C_{n+1}\ne0$ for all sufficiently large $n$, then $a_{n+1}=a_n^{2}-a_n+1$ for
all sufficiently large $n$.
\end{theorem}

\noindent Formalised as the
\lword{Erdos243/ReciprocalTailRigidity.lean}{1799}{sylvesterNext_eventually_of_centered_zero}{eventual
Sylvester recurrence}.  The proof is routine: take the larger of the two
thresholds and apply Theorem~\ref{res:step} at each later index.  Every later
section is directed at the hypothesis of this theorem, that is, at forcing $E$
to vanish eventually.

The defect identity has a second consequence, which is what makes a single
vanishing error worth having.

\begin{theorem}[zero is absorbing]\label{res:absorb}
Let $a,C,D:\N\to\N$ be an exact natural orbit, so
$C_{n+1}+D_n=a_nC_n$ and $D_{n+1}=a_nD_n$, and let
$E_n=\ctr(a_n,D_n,C_n)$.  Suppose the centring is strict, $|E_n|<C_n$ for
every $n$.  If $E_n=0$ then $E_{n+1}=0$.
\end{theorem}

\begin{proof}
Putting $E_n=0$ in Theorem~\ref{res:defect} gives
$\Delta_n C_{n+1}=-E_{n+1}$, so $C_{n+1}$ divides $E_{n+1}$.  A multiple of
$C_{n+1}$ of absolute value smaller than $C_{n+1}$ is zero, and
$|E_{n+1}|<C_{n+1}$ by hypothesis.
\end{proof}

\noindent Formalised as the
\lword{Erdos243/ReciprocalTailRigidity.lean}{2221}{centeredState_zero_absorbing}{absorption
of a vanishing centred state}, with the companion
\lword{Erdos243/ReciprocalTailRigidity.lean}{2248}{eventually_nonzero_of_zero_absorbing}{contrapositive
along a tail}: if zero is absorbing beyond some index and $E$ does not vanish
eventually, then $E$ is nowhere zero beyond that index.

Absorption changes the shape of the problem.  Either $E$ hits zero once and
stays there, or it is never zero.  Sections~\ref{sec:constant}
and~\ref{sec:periodic} address the special all-negative case only.  The
mixed-sign analysis
in Sections~\ref{sec:bounded} and~\ref{sec:mass} instead separates cofinally
negative errors from an eventually positive tail.

\section{Descent when the error is nonnegative}\label{sec:descent}

The hypothesis of the theorem below is the update law of
Proposition~\ref{res:update} read inside $\N$: the error is nonnegative at
every index, so the tail state never rises.  This is the easy case, and it is
settled in three lines.

\begin{theorem}[descent]\label{res:descent}
Let $C,E:\N\to\N$ satisfy $C_{n+1}+E_n=C_n$ for every $n$.  Then $E_n=0$ for
all sufficiently large $n$.
\end{theorem}

\begin{proof}
A one-line descent.  The relation forces $C_{n+1}\le C_n$, so $C$ is a nonincreasing sequence of
natural numbers and is eventually equal to its minimum; beyond that index
$C_{n+1}=C_n$, and the relation gives $E_n=0$.
\end{proof}

\noindent Formalised as the
\lword{Erdos243/ReciprocalTailRigidity.lean}{1837}{centeredState_eventually_zero}{stabilisation
of the nonnegative error}, on the
\lword{Erdos243/ReciprocalTailRigidity.lean}{1819}{antitone_nat_eventually_constant}{eventual
constancy of a nonincreasing sequence}.

The hypothesis is exactly one-sidedness: $C$ and $E$ take values in $\N$.  If
$E_n<0$ at infinitely many indices then $C$ rises there, nothing descends, and
the argument is unavailable.  That is the whole difficulty, and the rest of
this note is directed at it.
